% q0012.tex — Doubling the target (Industrial Organization / Cournot)
\question{
  id        = {0012},
  title     = {Doubling the Target},
  source    = {OBECON},
  year      = {2026},
  round     = {Seletiva IEO -- Phase A},
  language  = {English},
  topic     = {Micro},
  subtopic  = {Industrial Organization / Cournot Competition},
  type      = {objective},
  statement = {%
    Consider an infinite number of firms, each producing an identical good.
    Firm~$i$ has a marginal cost of production equal to $2^i$, for all
    $i \in \mathbb{Z}$. There are no fixed costs. The market demand is given
    by the function $Q = 2^{-P}$. A technological innovation is introduced
    that allows every firm to halve its marginal cost. Assuming firms operate
    under Cournot competition, how much will the joint profit of all the
    companies increase after the new technology?
  },
  choices   = {%
    \item $-50\%$
    \item $0\%$
    \item $50\%$
    \item $100\%$
    \item $200\%$
  },
  answer    = {B},
  solution  = {%
    Halving every firm's MC is equivalent to replacing $2^i$ with
    $2^{i-1} = 2^{(i-1)}$, i.e.\ shifting the index by $-1$. Since $i$
    ranges over all integers $\mathbb{Z}$, the set of marginal costs
    $\{2^i : i \in \mathbb{Z}\}$ is \emph{identical} before and after the
    shift (the set is unchanged; only the labeling shifts).

    Because the market structure --- the distribution of MCs across the
    infinite set of firms and the demand curve --- is identical before and
    after, the Cournot equilibrium (prices, quantities, profits) is
    \textbf{unchanged}.

    With infinitely many Cournot competitors whose MCs span $(0, \infty)$,
    the market approaches perfect competition with $P \to 0$, $\pi \to 0$
    regardless. The innovation leaves joint profit at 0 both before and after.

    Change in joint profit: $\mathbf{0\%}$. Answer: \textbf{(B)}.
  },
}
